\documentclass[reprint,amsmath,amssymb,floatfix,biblatex]{revtex4-2}

\usepackage{graphicx}% Include figure files
\usepackage{dcolumn}% Align table columns on decimal point
\usepackage{bm}% bold math
\usepackage{hyperref}% add hypertext capabilities
\usepackage{color}

%%
%%  bibtex stuff
%%
\usepackage[utf8]{inputenc}
\usepackage[english]{babel}


\newcommand{\bl}[1]{\color{blue}{#1}}

\begin{document}

\title{Exact Solutions of Burgers' Equation}% Force line breaks with \\
%\thanks{A footnote to the article title}%

\author{Ann Author}
\email{Ann.Author@students.plymouth.ac.uk}
\author{Second Author}%
 \email{Second.Author@students.plymouth.ac.uk}
 \author{Third Author}
  \email{Third.Author@students.plymouth.ac.uk}
\affiliation{School of Computing, Electronics and Mathematics, Plymouth University, Plymouth, PL4 8AA}

\date{\today}% It is always \today, today,
             %  but any date may be explicitly specified

\begin{abstract}
[This is where your brief abstract goes.] \bl{Find and discuss some exact analytical solutions to the viscous Burgers equation.}  
\end{abstract}

\maketitle

\section{Introduction}

This can provide some historical context, the state of the art and some motivation. Be reasonably general, but concise.

{\bl Instructions for your project choice are as follows:

\subsection{Oral Presentation}

The first part of this module is the preparation and delivery of an oral presentation on the topic stated in the title. 

You should address the following topics/problems related to the viscous Burgers' equation,
%
\begin{equation} \label{BURGERS}
 \partial_t u + u \partial_x u = \eta \partial_x^2 u \; , \; \eta > 0 \; .
\end{equation}
%

\begin{itemize}
\item Explain why this is a nonlinear diffusion equation. 
\item Show that (\ref{BURGERS}) has travelling wave solution of the form $u = v(x-ct)$.  
\item Map Burgers' equation to the heat equation using the Hopf-Cole transformation. Use this to find exact solutions of Burgers' equation.
\item Discuss the inviscid limit.  
\end{itemize}


\subsection{Scientific Paper}

In the second part you are to write a scientific article (or `paper) on your project topic. You must use this template typeset in \LaTeX. The paper will consist of a write-up of your oral presentation extended by at least one other application or natural extension of your topic, which was not covered by the oral presentation. The overall word count of your article must not exceed 4,000.}  

\section{Main Part I}

That's the first section of your main part. Formulae should be included as this is an applied maths paper, \textit{not} an essay. So please refrain from any ``subjective scientific writing''.  

An (enumerated) equation is included as follows,
%
\begin{equation}
  E = m c^2 \; .
\end{equation}
%
while two aligned equations would be:
%
\begin{eqnarray}
  (a + b)^2 &=& a^2 + 2 ab + b^2 \; , \\
  (a - b)^2 &=& a^2 - 2 ab + b^2 \; .
\end{eqnarray}
%
The format above (equations in extra lines) is called \textit{displaystyle}. If you write formulae within text, don't forget to enclose the maths expression in dollar symbols, e.g.\ like \verb|$a$|, which yields the maths letter $a$. This is proper maths  \textit{textstyle} --- I do not want to see just the ordinary letter ``a"! 

Some additional maths typesetting rules are:
%
\begin{itemize}
\item Generally avoid font sizes that are too small. In particular, in textstyle write fractions in the form $a/b$ (rather than the ugly $\frac{a}{b}$). 

\item Similar remarks apply to indices (powers, sub- and superscripts) in textstyle, in particular if they are more complicated expressions such as $x^{\frac{1}{1-x}}$.

\item Formulae should be treated as genuine parts of sentences, which includes the appropriate punctuation. 
\end{itemize}

If you want to include subsections, no problem:

%\newpage %enforces new column here; normally NOT required

\subsection{Subsection}

Always properly quote any work you are drawing from \cite{Weinstein:2008,Muller:2008,Einstein:1905ve}. This includes books, research articles, web sites and any other written resource you have been using. Anything else will be considered as plagiarism.

\section{Main Part II}

\begin{figure}[h]
\begin{center}
\includegraphics[scale=0.25]{./figures/Einstein1904.jpg}
\caption{\label{fig:Einstein}Einstein in 1904, at the time when he was working at the Swiss Patent Office \cite{Wiki:Einstein}.}
\end{center}
\end{figure}

Pictures and diagrams are as important as figures. If you use pdflatex (recommended) you can include pdf and jpg pictures as shown for Fig.~\ref{fig:Einstein}.
Note that figure files should be in the same directory as the tex file (unless you prescribe the explicit path before the filename).

\newpage

\section{Discussion and Conclusion}

Summarise and discuss your findings. Provide an outlook to further work.

\begin{acknowledgments}
Thank you!
\end{acknowledgments}

\bibliographystyle{plain} % We choose the "plain" reference style
\bibliography{refs} % Entries are in the "refs.bib" file

\end{document}

